\chapter{Introduction}

The title of Chapter 1 shall be Introduction. The introduction is certainly the most read chapter of any deliverable, and it largely determines the attitude of the reader/reviewer will have toward the work. Therefore, it is probably the most delicate part of the writing of a report.\\
It shall justify and highlight the problem posed, define the topic and explain the aim and scope of the work presented in the project report. It may also highlight the significant contributions from the investigation.\\
Also, it shall include the proposed solution and methodology used in brief.\\

\section{Motivation}

In this section, state the motivation w.r.t. project(brief about , what motivated you to take up this as a project) \\

\section{Objective of the project}

State the Objective precisely according to the suitability of the project.\\

\section{Organisation of the report}
This project report shall be presented in a number of chapters, starting with Introduction and ending with Summary and Conclusions. Each of the other chapters will have a precise title reflecting the contents of the chapter. A chapter can be subdivided into sections, subsections and sub subsection so as to present the content discretely and with due emphasis. When the work comprises two or more mutually independent investigations, the project report may be divided into two or more parts, each with an appropriate title. However, the numbering of chapters will be continuous right through. In this section make sure to mention in what chapter what is explained.